% 天体（综述）
% license CCBYSA3
% type Wiki

本文根据 CC-BY-SA 协议转载翻译自维基百科\href{https://en.wikipedia.org/wiki/Astronomical_object}{相关文章}。

一个天文对象、天体、恒星对象或天界对象，是指在可观测宇宙中自然存在的物理实体、关联或结构。在天文学中，“对象”（object）和“天体”（body）这两个术语常被交替使用。然而，“天文天体”“天体”或“天界天体”特指单一的、紧密结合的、连续的物理体；而“天文对象”或“天体对象”则可以指更为复杂、结合程度较弱的结构，其中可能包含多个天体，甚至包含带有次级结构的其他对象。

天文对象的例子包括行星系统、星团、星云和星系；而小行星、卫星（moons）、行星和恒星则属于天文天体。一颗彗星可以同时被视为天体和对象：当指其由冰与尘埃组成的冻结核时，它是一个天体；而当描述其整体，包括弥散的彗发与彗尾时，它则属于天文对象。
\subsection{历史}
根据美国宇航局（NASA）的天体物理学家推测，最早的天文对象大约在 136 亿年前出现，大约是在大爆炸之后 2 亿年。随着时间推移，光子通过引力聚集，形成了最初的恒星和星系。

恒星、行星、星云、小行星与彗星等天文对象已经被人类观测了数千年，尽管早期文明多将这些天体视为神祇。早期文化认为天体的运动极为重要，因为人们依赖这些天体在长距离航行中导航、区分季节，以及判断农作物的播种时机。在中世纪，各文明开始更系统地研究这些天体的运动。中东地区的数位天文学家对恒星和星云进行了详细描述，并依据这些天体的运动制定更准确的历法。在欧洲，天文学家更多地专注于制造用于观测天体的仪器，以及编写教材、指南与设立大学，以系统教授天文学知识。

在科学革命期间，1543 年，尼古拉·哥白尼发表了日心模型。该模型描述地球与所有其他行星一样，作为天文天体围绕太阳运行，而太阳位于太阳系的中心。约翰内斯·开普勒发现了行星运动三大定律，这些定律刻画了天体轨道的特性，并用于进一步改进日心模型。1584 年，焦尔达诺·布鲁诺提出遥远恒星即为其他太阳，这是数百年来第一次有人提出这一观点。伽利略·伽利莱是最早使用望远镜观测天空的天文学家之一。1610 年，他观测到了木星的四颗最大卫星（现称伽利略卫星）。伽利略还观测了金星的相位、月球上的环形山，以及太阳上的黑子。天文学家埃德蒙·哈雷精确预言了 1758 年后来以他名字命名的哈雷彗星的回归。1781 年，威廉·赫歇尔爵士发现了天王星，这是第一颗肉眼不可见但经望远镜发现的行星。

在 19 世纪与 20 世纪，新技术与科学创新使人类得以大幅拓展对天文学与天文对象的理解。大型望远镜与天文台开始建造，科学家们开始利用照相底板记录月球与其他天体的影像。人眼不可见的新波段被发现，并制造出能在不同电磁波段观测天体的新型望远镜。约瑟夫·冯·夫琅和费与安杰洛·塞基开创了光谱学，使他们得以分析恒星与星云的组成。许多天文学家也得以依据双星系统的轨道参数推算其质量。计算机开始用于处理大量恒星观测数据，而诸如光电光度计等新技术让天文学家能够精确测量恒星的颜色与光度，从而推算其温度与质量。
1913 年，天文学家埃伊纳尔·赫茨普龙与亨利·诺里斯·罗素分别独立发展了赫茨普龙–罗素图，它根据恒星的颜色与光度进行绘制，使恒星性质的研究变得更为系统。人们发现恒星普遍分布在图中的一个带状区域，即主序星。1943 年，威廉·威尔逊·摩根与菲利普·基尔兹·基南在赫茨普龙–罗素图基础上制定了更精细的恒星分类体系。天文学家们也开始讨论银河系以外是否存在其他星系。这一讨论最终在埃德温·哈勃将仙女座星云确认为另一个星系（以及许多距离银河系更远的星系）后尘埃落定。
\subsection{星系及更大尺度}
宇宙可以被视为具有层级结构。在最大的尺度上，其基本构成单位是星系。星系会进一步聚集成星系群与星系团，并常常位于更大的超星系团之中。这些超星系团沿着巨大的丝状结构排列，而丝状结构之间则是近乎空旷的巨大空洞，整体形成一张贯穿可观测宇宙的宇宙网。

星系具有多种形态，包括不规则星系、椭圆星系与盘状星系。这些形态取决于其形成与演化历史，包括与其他星系的相互作用，这些作用有时会导致星系并合。盘状星系包含透镜状星系与螺旋星系，并可能具有螺旋臂与明显的星系晕等结构。大多数星系核心处包含一个超大质量黑洞，这可能导致活动星系核的形成。星系还可能拥有卫星，例如矮星系与球状星团。
\subsection{星系内部}
星系内部的各种组成成分由气体物质在引力自吸引的层级聚集过程中形成。在这一层级中，最基本的结构单元是恒星，它们通常由不同的致密星云聚集而成的星团中形成。恒星多样性的根本原因在于其质量、组成与演化阶段的差异。恒星可以组成多星系统，以层级方式互相绕转。

围绕新形成恒星的原行星盘通过吸积过程能够形成行星系统以及各种小型天体，例如小行星、彗星与碎片等。

恒星的不同类型由赫茨普龙–罗素图（H–R 图）展示，该图将恒星的绝对光度与表面温度对应绘制。每一颗恒星在该图上都遵循一条演化轨迹。如果这条轨迹穿过某些本征变星的分布区域，那么恒星的物理性质就可能促使其成为变星。例如，不稳定带是 H–R 图上的一个区域，其中包括天琴座δ型变星、RR Lyrae 型变星与造父变星等。

随着演化的继续，恒星可能抛出部分大气，形成星云。这一过程可能较为平稳，从而形成行星状星云；也可能以超新星爆发的方式发生，并留下超新星遗迹。根据恒星的初始质量以及是否存在伴星，其生命末期可能以致密天体的形式存在：白矮星、中子星或黑洞。
\subsection{形状}
\begin{figure}[ht]
\centering
\includegraphics[width=8cm]{./figures/1288a73312132d08.png}
\caption{} \label{fig_TianTI_1}
\end{figure}
